% ============================================================
% PREFAZIONE (CAPITOLO NON NUMERATO)
% ============================================================
\chapter*{Istruzioni per l'uso}

\markboth{Prefazione}{}

Quello che hai davanti non è un manuale di armonia nel senso tradizionale del termine. Non è un testo di consultazione, né pretende di essere esaustivo. Queste pagine sono, molto più semplicemente, i miei appunti sparsi: il tentativo di mettere ordine nel meraviglioso e a volte caotico mondo della sintassi tonale, così come l'ho appreso io stesso.

Da studente di musica, ma con un passato da fisico e un presente da insegnante di matematica, ho affrontato questa disciplina con una certa "deformazione professionale". Ho cercato schemi, logiche e principi primi, convinto che l'armonia, prima di essere arte, sia una scienza di straordinaria coerenza. Questi appunti, quindi, non seguono un ordine tematico o enciclopedico, ma un percorso di apprendimento progressivo, pensato per costruire la conoscenza un mattone alla volta.

Gli appunti qui trascritti derivano dalle lezioni del \textbf{prof. M° Davide Vallini}, dalla manualistica classica e, lo ammetto, dal pesante contributo dei Grandi Modelli di Linguaggio (LLM)\footnote{In particolare Google Gemini e Claude di Anthropic.}. Proprio per questo, non troverai qui un eserciziario: per quello, ti rimando con convinzione ai manuali "ufficiali" citati in bibliografia. Loro sono le fonti autorevoli; questi sono solo gli appunti di un compagno di viaggio.

Per navigare in queste pagine, ho trovato utile una struttura che tradisce la mia formazione. Ogni concetto è affrontato seguendo tre passi:
Per navigare in queste pagine, ho trovato utile una struttura che tradisce la mia formazione. Ogni concetto è affrontato seguendo questi passi:
\begin{itemize}
	\item \textbf{Definizione.} Come in matematica, ci serve un linguaggio comune per sapere esattamente di cosa stiamo parlando. Le definizioni non si discutono, si imparano.
	\item \textbf{Regola.} Una volta compreso il concetto, la regola ci dice come applicarlo nella pratica della scrittura armonica.
	\item \textbf{Esempio.} Qui vediamo la regola in azione.
	\item \textbf{Errore.} Non un divieto fine a se stesso, ma l'analisi del \textit{perché} qualcosa non funziona. Capire la logica di un errore è il primo passo per evitarlo, e un giorno, forse, per infrangerlo con consapevolezza.
\end{itemize}

A questa struttura si aggiungono, in alcuni punti del manuale, degli approfondimenti racchiusi in box dalla cornice squadrata e dallo sfondo grigio chiaro, intitolati \textbf{"Con la fisica"} e \textbf{"Con la matematica"}. Questi non sono essenziali per lo studio dell'armonia in senso stretto, ma offrono una prospettiva complementare a chi, come me, trova che la comprensione e l'apprezzamento di un fenomeno passi anche attraverso la sua formalizzazione quantitativa mediante le equazioni, i numeri e le formule.

Perché questa scelta? Perché l'armonia tonale non è solo una convenzione culturale o un insieme di regole stilistiche tramandate per tradizione. Essa è profondamente radicata nelle leggi della fisica (la serie degli armonici, i rapporti di frequenza, i battimenti) e può essere descritta con straordinaria precisione attraverso gli strumenti della matematica (l'aritmetica modulare, la teoria degli insiemi, le funzioni ricorsive). Comprendere \textit{perché} una quinta giusta è consonante, \textit{perché} il tritono genera tensione, o \textit{perché} le progressioni armoniche seguono pattern prevedibili, non toglie nulla alla bellezza della musica. Al contrario, ne rivela la natura di linguaggio universale, ponte tra arte e scienza, tra intuizione estetica e rigore logico.

Se questi box ti sembrano superflui o troppo tecnici, saltali pure senza remore: il percorso principale procede in modo del tutto autonomo. Ma se, come me, trovi che "sapere il perché" arricchisca il "sapere il come", allora potresti trovarvi qualche spunto interessante per guardare la musica con occhi diversi.

Ora, lasciami essere brutalmente onesto: leggere questi appunti non ti insegnerà l'armonia. L'armonia non si impara leggendo, si impara scrivendo. Svolgere esercizi è l'unica via. Falli. E, ancora più importante, \textbf{falli correggere dal tuo insegnante}. Un errore che non viene visto è un'occasione di apprendimento sprecata. L'occhio esterno di un maestro vede le incertezze che a noi sfuggono e trasforma uno scarabocchio in una lezione.

Questa centralità dell'errore non è una mia fissazione, ma la convinzione che sia il vero motore della conoscenza. Nessuno lo ha espresso meglio del filosofo Karl Popper:

\begin{quote}
\emph{	Evitare errori è un ideale meschino: se non osiamo affrontare problemi che siano così difficili da rendere l'errore quasi inevitabile, non vi sarà allora sviluppo della conoscenza. In effetti, è dalle nostre teorie più ardite, incluse quelle che sono erronee, che noi impariamo di più.} (K. Popper, 1972)
\end{quote}

Quindi, ti auguro un buono studio, un buon lavoro e, soprattutto, buoni errori.
